% 出席確認クイズ 第14回　最終発表：AI活用提案
%   attendance/attendance14.html から生成（解説は本ガイド用に加筆）。

\quizq{1}{AI活用提案で最初に明確にすべきことはどれでしょうか?}%
  {使うAIツールの名前}%
  {\correct{解決したい課題と目的}}%
  {予算だけ}%
  {発表スライドの色}%
  {正解は (b)。提案は「何のために」から始めます。解決したい課題と目的が明確でないと、AIを使うこと自体が目的化します（第7回の失敗要因）。(a) ツールの選定は目的が決まった後、(c) 予算と (d) スライドの色は本質ではありません。}

\quizq{2}{提案の実現可能性を検討する際の観点として適切なものはどれでしょうか?}%
  {見た目の派手さ}%
  {流行しているかどうかだけ}%
  {\correct{データの有無、コスト、組織の受け入れ、法規制}}%
  {発表時間の長さ}%
  {正解は (c)。実現可能性は、必要なデータの有無、コスト、組織の受け入れ、法規制などの観点から検討します。(a)(b)(d) は実現可能性と無関係です。この4観点は第8・12・6・13回で扱った内容に対応しています。}

\quizq{3}{提案の効果を示す方法として適切なものはどれでしょうか?}%
  {感想だけを書く}%
  {効果は示さなくてよい}%
  {\correct{導入前後で比較できる指標(KPI)を設定する}}%
  {導入費用の高さで示す}%
  {正解は (c)。効果は、導入前後で比較できる具体的な指標（KPI）を設定して示します（第7回）。(a) 感想や (d) 費用の高さでは効果を示せず、(b) 効果を示さない提案は採用されません。}

\quizq{4}{AIを使って作成した提案を発表するときの姿勢として適切なものはどれでしょうか?}%
  {自分の判断は不要である}%
  {すべてAIの出力のまま提出する}%
  {\correct{AIの寄与と自分の判断・検証を区別して示す}}%
  {AIを使ったことは隠す}%
  {正解は (c)。本講義の AI Policy のとおり、AI がどこに寄与し、自分が何を判断・検証したかを区別して示します。(a)(b) は自分の判断の放棄、(d) は透明性の否定です。}

\quizq{5}{相互評価(ピアレビュー)の目的として適切なものはどれでしょうか?}%
  {順位を付けることだけ}%
  {\correct{他者の視点で提案の強みと改善点を見つける}}%
  {形式的に済ませること}%
  {発表者を批判すること}%
  {正解は (b)。相互評価は、他者の視点から強みと改善点を見つけ、提案をよりよくするために行います。(a) 順位付けや (d) 批判が目的ではなく、(c) 形式的に済ませては意味がありません。評価する側も、提案を評価する観点を身につける学習機会になります。}
